\scanpage{262}
\mspage{3}

Quand on parlera des ``composantes connexes'' de $E$, il s'agira toujours des classes suivant cette relation d'équivalence.

Nous \unsure{identifions|identifierons} l'ens des flèches de transition \inserted{$\mathrm{Fl.tr.}(I)$} de $I$ à l'intervalle d'entiers $[1,n]$, donc à une partie de $\mathbb{Z}$. Si on a une partie $E\subset\mathrm{Fl.tr.}(I)$, \deleted{\unsure{elles sont}} \inserted{on y a} donc \deleted{\unsure{aussi}} la notion de flèches de transition ``connectées'' dans $E$, et la notion de connexité pour $E$.

Si $i<j$, $\underline{u}_i$ et $\underline{u}_j$ sont connectées \unsure{sss|ssi} \unsure{toutes} les flèches de transition intermédiaires entre $\underline{u}_i$ et $\underline{u}_j$ sont dans $E$. Et $E$ est connexe sss il est soit vide, soit formé d'une suite de flèches de transition dont chacune est successeur immédiat de celle qui \unsure{la|le} précède

\[E=\{\underline{u}_i,u_{i+1},\ldots u_{i+l-1}\},\]

où $l$ est le \unsure{``}cardinal de $E$. Une telle $E$ est appelée une \uline{\unsure{chaine|chaîne}} de flèches de transition, et $l$ est sa longueur.

Une chaîne est dite \uline{composable} si \unsure{le|la} \deleted{flèche $\unsure{u}_{i+l-1}$ \illegible{1}} \deleted{\unsure{suite}} \inserted{composé $\underline{u}_{i+l-1}\circ\cdots\underline{u}_{i+1}\circ\underline{u}_i$ a un sens, i.e. si elles sont toutes de $\mathord{=}$ sens, soit direct, soit rétrograde.} On désigne par $\mathrm{Fl.tr.}^{+}(I)$, $\mathrm{Fl.tr.}^{-}(I)$ \deleted{les deux} \inserted{l'ens.} des fl. de tr. directes, ou rétrogrades, respectivement.

Leurs \deleted{\unsure{con}} composantes connexes ne sont autres que les chaînes composables maximales. \deleted{\unsure{qu'on appellera}} \inserted{Ainsi, les} \deleted{les \uline{chaînes} \uline{fondamentales}} chaînes composables maximales forment une partition de l'ens. des \deleted{chaînes} \inserted{flèches de} \unsure{et|de} transition, en intervalles d'entiers mutuellement disjoints. Ces chaînes forment elles-$\mathord{=}$s un ens. totalement ordonné

\[C_1,C_2,\ldots,C_r\qquad(1\leq r\leq n)\]

L'entier $r$ s'appelle la ``\uline{longueur utile}'' de l'intervalle

% archive sheet number: 261
